\section{Conclusion}
To conclude, it is worth highlighting the successful creation of a scalable and modular benchmark for evaluating the PDDL planning capabilities of Large Language Models. The framework was designed to be extensible to a wide range of problems and models, with a strong emphasis on producing easily interpretable and traceable results. The \textit{raw, parsed, scored} output structure ensures diagnostic precision, allowing for a clear separation between generation, parsing, and validation failures.

\subsection{General Conclusions on LLM Planning Capabilities}
Our evaluation of several models highlighted critical insights about the current state of LLM-based planning: the main finding is that while LLMs often handle PDDL syntax and action vocabulary, their ability to perform robust, long-horizon, and goal-directed reasoning is still fragile.
Models frequently generate action sequences that are syntactically correct and locally executable, as pointed out by generally high \texttt{Executability Ratios} and low \texttt{Hallucination Rates}, but these plans often fail to achieve the global goal, leading to low overall \texttt{Success Rates}. The significant gap between first-attempt success (\texttt{FASR}) and the final success rate (\texttt{SR}) highlights a strong dependency on iterative repair; this suggests that current models often behave more like retry-dependent generators than deterministic planners.

\subsection{Limitations of the Study}
It is also worth noting the limitations of this study: the evaluation is constrained by the available hardware, since these LLM models require significant resources to run effectively.
The observed performance is also tied to the specific prompts and iterative repair strategies employed; different approaches could yield different results.

\subsection{Hints for Future Work}
The modularity of the benchmark creates several possibilities for future work and improvement:
\begin{itemize}
    \item \textbf{Single Domain Testing:} The set of PDDL problems could be reduced to a single reasoning type and the difficulty increased gradually to see where a given LLM starts to fail at planning.
    
    \item \textbf{Advanced Interaction Protocols:} Future work could explore more sophisticated interaction protocols. This includes implementing more informative feedback mechanisms for the iterative repair loop, providing models with specific details about precondition failures (e.g., "action `(A)` failed because precondition `(P)` was false") to better assess their self-correction abilities.
    
    \item \textbf{Broader Model Evaluation:} The benchmark should be used to evaluate a broader and more diverse set of models, to build a more comprehensive landscape of LLM planning skills.

    \item \textbf{Single Model Multi-Release Evaluation:} The benchmark could be used to evaluate different releases of the same model, to test how its planning capabilities evolve across versions.
    
    \item \textbf{Deeper Analytical Insights:} The analytical layer can be extended to perform deeper correlational studies, such as linking specific model capability profiles to performance on domains with certain structural properties (e.g., high branching factor or sparse rewards), which could reveal fundamental strengths and weaknesses of different architectures.

    \item dal reasoning text per essere sicuri di avere il reasoning corretto si puo implementare una ia per trovare il piano dal reasoning invece che un approccio hard coded (come quello implementato qui)
\end{itemize}
